% This file recreates the Knucklebones puzzle page. The diagram is built from simple
% LaTeX shapes so the proof remains complete even without the original board image.
\begin{multicols}{2}
\DndrSectionBanner{5.0 The Puzzle Minigame Vault}
When players choose a puzzle resolution path, the Architect operates as the obstacle and uses the standard cooperative game format.

\subsection*{Knucklebones}
Knucklebones resolves Puzzle Quests. It can represent picking a lock, disabling a magical trap, unraveling an ancient ritual, or overcoming any obstacle requiring wit rather than force.

The Architect plays against the party. If the party wins every active Knucklebones board, the Quest succeeds. If the Architect wins even one board, the Quest fails.

\subsection*{Setup}
Each side receives a board of 3 columns with 3 spaces each. The Architect starts with a 6 in the center square. All participants use standard d6s.

For groups of 4 or more players, the Architect assigns a helper and runs a second board simultaneously. The party must win both boards.

Before play, cover Class Powers equal to the current Challenge Rating. Covered powers cannot be used during the puzzle.

\subsection*{Taking Turns}
The Architect takes the first turn. On your turn:
\begin{enumerate}
  \item Roll a d6.
  \item Place the die into any column on your board.
  \item Optionally use an Item or Class Ability.
\end{enumerate}
Any player may spend Applause Tokens, activate Items, or use available Class Abilities to assist any party board.

\DndrCallout{Keep the Puzzle Alive}{Describe the puzzle as the board changes. A removed die can be a lock tumbler clicking into place, a rune fading, or part of a ritual circle collapsing.}

\begin{center}
\begin{tikzpicture}[x=0.42in,y=0.42in]
  \foreach \x in {0,1,2}{\foreach \y in {0,1,2}{\draw[fill=DndrBlueLight,draw=DndrBlue] (\x,\y) rectangle +(1,1);}}
  \node[fill=DndrBlue,text=white,font=\bfseries] at (1.5,1.5) {4};
  \node[fill=DndrBlue,text=white,font=\bfseries] at (1.5,0.5) {4};
  \node[fill=DndrBlue,text=white,font=\bfseries] at (1.5,2.5) {4};
  \node[rotate=90,font=\bfseries\scriptsize,DndrBlue] at (-0.25,1.5) {KNUCKLEBONES};
\end{tikzpicture}
\end{center}

\subsection*{Scoring}
When the game ends, each die is worth its face value. Matching dice in the same column multiply their value:
\begin{itemize}
  \item Two matching dice are worth double each die's value.
  \item Three matching dice are worth triple each die's value.
\end{itemize}
A column containing 5, 5 is worth 20 points. A column containing 4, 4, 4 is worth 24 points.

\subsection*{Removing Opponent Dice}
When you place a die into a column, remove one matching die from the same column on your opponent's board.

\DndrCallout{Knucklebones Tip}{Removing an opponent's high-value dice is often more powerful than adding points to your own board.}

\subsection*{Winning the Puzzle}
The game ends when either side fills all 9 spaces. Calculate both scores. Higher score wins; ties favor the Architect. The party must win every active board. Resolve success, failure, and rewards as described in Section 4.2.
\end{multicols}

